\documentclass[11pt,a4paper]{article}

\usepackage[T1]{fontenc}
\usepackage[utf8]{inputenc}
\usepackage{lmodern}
\usepackage[ngerman]{babel}
\usepackage[a4paper,margin=2.2cm]{geometry}
\usepackage{array}
\usepackage{longtable}
\usepackage{tabularx}
\usepackage{xcolor}
\usepackage{hyperref}

\setlength{\parindent}{0pt}
\setlength{\parskip}{0.5em}
\renewcommand{\arraystretch}{1.35}

\newcommand{\answerline}[1]{\rule{#1}{0.4pt}}
\newcommand{\term}[1]{\fcolorbox{black!20}{blue!6}{\strut #1}}

\title{\textbf{Aufgabenblatt: CQRS und Event Sourcing}}
\author{Modul 321 -- Lektion 03}
\date{}

\begin{document}

\maketitle

\textbf{Name:} \answerline{7cm} \hfill \textbf{Datum:} \answerline{3.5cm}

\vspace{0.8em}

\fcolorbox{blue!40}{blue!6}{
  \parbox{\dimexpr\linewidth-2\fboxsep-2\fboxrule-6pt\relax}{
    \textbf{Ziel der Lektion:} Sie können CQRS und Event Sourcing fachlich sauber
    erklären, mit Beispielen anwenden und deren Nutzen sowie Grenzen in
    verteilten Systemen begründet einordnen.
  }
}

\section*{Hinweise}
\begin{itemize}
  \item Antworten Sie in ganzen Sätzen oder präzisen Stichworten.
  \item Begründen Sie fachliche Entscheidungen knapp, aber nachvollziehbar.
  \item Kleine Skizzen, Event-Flows oder Tabellen sind ausdrücklich erlaubt.
\end{itemize}

\section*{1. Begriffe sauber erklären}
Erklären Sie die folgenden Begriffe in \textbf{1 bis 3 eigenen Sätzen}.

\begin{longtable}{|>{\raggedright\arraybackslash}p{0.28\linewidth}|>{\raggedright\arraybackslash}p{0.64\linewidth}|}
\hline
\textbf{Begriff} & \textbf{Ihre Erklärung} \\
\hline
\term{CQRS} & \\[2.4cm] \hline
\term{Command} & \\[2.0cm] \hline
\term{Query} & \\[2.0cm] \hline
\term{Event Sourcing} & \\[2.4cm] \hline
\end{longtable}

\section*{2. Command oder Query?}
Ordnen Sie die folgenden Beispiele zu. Schreiben Sie jeweils \textbf{Command} oder \textbf{Query} hin und begründen Sie kurz.

\begin{tabularx}{\linewidth}{|>{\raggedright\arraybackslash}p{0.44\linewidth}|>{\centering\arraybackslash}p{0.16\linewidth}|>{\raggedright\arraybackslash}X|}
\hline
\textbf{Beispiel} & \textbf{Typ} & \textbf{Begründung} \\
\hline
\texttt{CreateBooking} & & \\[0.9cm] \hline
\texttt{GetOpenBookings} & & \\[0.9cm] \hline
\texttt{CancelInvoice} & & \\[0.9cm] \hline
\texttt{GetCourseOverview} & & \\[0.9cm] \hline
\end{tabularx}

\section*{3. Warum CQRS?}
Eine Plattform für Kursanmeldungen hat folgende Anforderungen:
\begin{itemize}
  \item Viele Lernende rufen ständig Übersichten und freie Plätze ab.
  \item Schreibvorgänge sind seltener, müssen aber Regeln strikt prüfen.
  \item Die Startseite zeigt mehrere vorberechnete Statistiken.
\end{itemize}

Bearbeiten Sie die Fragen:
\begin{enumerate}
  \item Warum könnte ein gemeinsames CRUD-Modell hier unpraktisch werden?
  \item Welchen Vorteil bringt ein getrenntes \term{Write-Model}?
  \item Welchen Vorteil bringt ein getrenntes \term{Read-Model}?
  \item Weshalb darf das Read-Model Daten anders strukturieren als das Write-Model?
\end{enumerate}

\section*{4. Event Sourcing nachvollziehen}
Betrachten Sie folgenden Event-Stream eines Bankkontos:

\begin{quote}
\texttt{AccountOpened(0)} \\
\texttt{MoneyDeposited(150)} \\
\texttt{MoneyWithdrawn(40)} \\
\texttt{MoneyDeposited(25)}
\end{quote}

Bearbeiten Sie die Aufgaben:
\begin{enumerate}
  \item Berechnen Sie den aktuellen Kontostand.
  \item Beschreiben Sie in eigenen Worten, was beim \term{Replay} dieses Streams passiert.
  \item Welcher Vorteil entsteht, wenn diese Ereignisse dauerhaft gespeichert bleiben?
  \item Nennen Sie einen möglichen Nachteil dieses Ansatzes.
\end{enumerate}

\section*{5. Command oder Event?}
Ordnen Sie die folgenden Begriffe zu und erklären Sie jeweils kurz, warum:

\term{Command}, \term{Event}

\begin{enumerate}
  \item \texttt{PlaceOrder}
  \item \texttt{OrderPlaced}
  \item \texttt{ReserveSeat}
  \item \texttt{SeatReserved}
\end{enumerate}

\vspace{0.8em}
\textbf{Zusatzfrage:} Warum kann ein System einen Command ablehnen, ein Event aber nicht einfach wieder "wünschen"?

\section*{6. CQRS und Event Sourcing kritisch beurteilen}
Beantworten Sie die Aussagen mit \textbf{richtig / falsch} und begründen Sie kurz.

\begin{enumerate}
  \item CQRS bedeutet zwingend zwei getrennte Anwendungen.
  \item Ein Read-Model darf aus Events neu aufgebaut werden.
  \item Event Sourcing ist immer die beste Wahl für einfache CRUD-Formulare.
  \item Commands und Events beschreiben fachlich dasselbe aus derselben Perspektive.
\end{enumerate}

\section*{7. Read-Model entwerfen}
Ein Ticketsystem soll folgende Anzeige liefern:
\begin{itemize}
  \item Veranstaltungstitel
  \item Anzahl freie Plätze
  \item Anzahl reservierte Tickets
  \item Anzahl bezahlte Tickets
\end{itemize}

Die Write-Seite arbeitet mit Commands wie \texttt{ReserveTicket} und \texttt{PayTicket}.

Bearbeiten Sie die Aufgaben:
\begin{enumerate}
  \item Beschreiben Sie ein mögliches Read-Model für diese Ansicht.
  \item Welche Events könnten dieses Read-Model aktualisieren?
  \item Warum wäre dieses Read-Model für die Anzeige nützlicher als rohe Event-Daten?
\end{enumerate}

\section*{8. Grössere Praxisaufgabe}
Sie entwerfen eine Plattform zur \textbf{Reservation von Gruppenräumen} für eine Schule. Die Plattform soll:
\begin{itemize}
  \item Räume buchbar machen,
  \item freie Kapazitäten schnell anzeigen,
  \item Stornierungen unterstützen,
  \item eine Historie aller Änderungen behalten,
  \item ein Dashboard mit Auslastung pro Tag anzeigen.
\end{itemize}

Bearbeiten Sie dazu:
\begin{enumerate}
  \item Beschreiben oder zeichnen Sie eine Lösung mit \term{CQRS}.
  \item Definieren Sie mindestens \textbf{vier Commands}.
  \item Definieren Sie mindestens \textbf{vier Events}.
  \item Beschreiben Sie ein mögliches \term{Write-Model} oder Aggregat mit wichtigen Fachregeln.
  \item Beschreiben Sie mindestens \textbf{zwei Read-Modelle} für unterschiedliche Ansichten.
  \item Erklären Sie, wie \term{Event Sourcing} in diesem Szenario die Historie und den Neuaufbau von Ansichten unterstützen würde.
  \item Nennen Sie \textbf{zwei zusätzliche technische Herausforderungen}, die bei dieser Lösung entstehen.
\end{enumerate}

\section*{9. Umsetzungsnahe Zusatzaufgabe}
Formulieren Sie für die Plattform aus Aufgabe 8 einen möglichen Ablauf in \textbf{5 bis 8 Schritten}, wie eine Raumreservation verarbeitet wird.

Verwenden Sie dabei möglichst die Begriffe:
\begin{itemize}
  \item Command
  \item Validierung
  \item Event
  \item Projection / Read-Model
  \item eventual consistency
\end{itemize}

\section*{Abgabe / Besprechung}
Die Aufgaben werden in der Lektion besprochen. Halten Sie Ihre Antworten so fest,
dass Sie diese fachlich sauber vorstellen und begründen können.

\end{document}
